\subsection{有序对公理}

正如我们在§1第1节中所说，符号 \(\supset\) 是集合论中权重为 2 的实体化符号。若 T、U 是项，则 \(\supset\) TU 是一个项，在实际中记作 (T, U)。在此记号下，有序对公理如下：

\[
(\forall x) (\forall x ^ {\prime}) (\forall y) (\forall y ^ {\prime}) (((x, y) = (x ^ {\prime}, y ^ {\prime})) \Longrightarrow (x = x ^ {\prime} \quad \text { 且 } \quad y = y ^ {\prime})). \tag {A3.}
\]

由于关系“ \(x = x'\) 并且 \(y = y'\)'' 蕴含 \((x, y) = (x', y')\) 由C44（第一章，§5，第2节），关系 \((x, y) = (x', y')\) 等价于“ \(x = x'\) 并且 \(y = y'\)''.

¶ 关系 \((\exists x)(\exists y)(z=(x,y))\) 写作“z是有序对”。若z是有序对，则关系

\[
(\exists y) (z = (x, y)), \quad (\exists x) (z = (x, y))
\]

关于 \(x\) 并且 \(y\) 分别地；这是A3的直接推论。

这些项

\[
\tau_ {x} ((\exists y) (z = (x, y))) \quad \text { 且 } \quad \tau_ {y} ((\exists x) (z = (x, y)))
\]

分别记为 \(\mathbf{pr}_1z\) 和 \(\mathbf{pr}_2z\) ，并称为 \(z\) 的第一坐标（或第一投影）和第二坐标（或第二投影）。如果 \(z\) 是有序对，则关系 \((\exists y)(z = (x, y))\) 等价于 \(x = \mathbf{pr}_1z\) ，而关系 \((\exists x)(z = (x, y))\) 等价于 \(y = \mathbf{pr}_2z\) （第一章，第5节，第3条）。

关系 \(z = (x, y)\) 等价于“ \(z\) 是一个有序对，并且 \(x = \text{pr}_1z\) 并且 \(y = \text{pr}_2z\) ''；因为后一个关系等价于

\[
(\exists x ^ {\prime}) (\exists y ^ {\prime}) (\exists x ^ {\prime \prime}) (\exists y ^ {\prime \prime}) (z = (x ^ {\prime}, y ^ {\prime}) \text {   且   } z = (x, y ^ {\prime \prime}) \text {   且   } z = (x ^ {\prime \prime}, y));
\]

由A3，“ \(z = (x', y')\) 并且 \(z = (x, y'')\) 并且 \(z = (x'', y)\)”等价于“ \(z = (x, y)\) 并且 \(x = x'\) 并且 \(x = x''\) 并且 \(y = y'\) 并且 \(y = y''\) ”；因此由C33（第一章，§4，第3节），“ \(z\) 是一个有序对，并且 \(x = \text{pr}_1 z\) 并且 \(y = \text{pr}_2 z\) ”等价于

\[
z = (x, y) \text { 且 } (\exists x ^ {\prime}) (\exists y ^ {\prime}) (\exists x ^ {\prime \prime}) (\exists y ^ {\prime \prime}) (x = x ^ {\prime} \text { 且 } x = x ^ {\prime \prime} \text { 且 } y = y ^ {\prime} \text { 且 } y = y ^ {\prime \prime}),
\]

这证明了我们的断言。显然我们有

\[
\operatorname{pr} _ {1} (x, y) = x, \quad \operatorname{pr} _ {2} (x, y) = y,
\]

且关系 \(z = (\mathrm{pr}_1(z), \mathrm{pr}_2(z))\) 等价于“ \(z\) 是一个有序对”。

设 \(R\{x, y\}\) 为一个关系，其中字母x和y互不相同且出现在R中。设z为一个字母，与x和y均不同，且不出现在R中。设 \(S\{z\}\) 表示该关系

\[
(\exists x) (\exists y) (z = (x, y) \text {   且   } R \{x, y \});
\]

这是一个比R少一个字母的关系，并且它等价于“z是一个有序对且 \(R\{pr_{1}z, pr_{2}z\}\) ”，这由以下事实得出 \(z = (x, y)\) 等价于“z 是一个有序对且 \(x = pr_{1}x\) 并且 \(y = pr_{2}z\) ''，并且根据准则C33（第一章，§4，第3节）

以及C47（第一章，§5，第3节）。由此立即得出 \(R\{x, y\}\) 等价于 \(S\{(x, y)\}\) ，并且也

\[
(\exists z) (z = (x, y) \text {   且   } S \{z \})
\]

由C47。

这意味着我们可以将对象x和y之间的关系解释为由这两个对象构成的有序对所具备的一种性质。

